\documentclass[11pt]{article}
\usepackage[T1]{fontenc}
\usepackage{lmodern,microtype}
\usepackage[margin=1in]{geometry}
\usepackage{amsmath,amssymb,amsthm,mathtools,bm}
\usepackage{booktabs,array,longtable,enumitem}
\usepackage[dvipsnames]{xcolor}
\usepackage{xurl}
\usepackage[colorlinks=true,linkcolor=MidnightBlue,citecolor=MidnightBlue,urlcolor=MidnightBlue]{hyperref}
\usepackage{fancyhdr}
\pagestyle{fancy}
\fancyhf{}
\fancyhead[L]{\small Binary broadcasts: analytic renormalization}
\fancyhead[R]{\small ProveIt research note}
\fancyfoot[C]{\thepage}
\setlength{\headheight}{14pt}
\setlength{\emergencystretch}{2em}
\newtheorem{theorem}{Theorem}[section]
\newtheorem{proposition}[theorem]{Proposition}
\newtheorem{lemma}[theorem]{Lemma}
\newtheorem{corollary}[theorem]{Corollary}
\theoremstyle{definition}
\newtheorem{definition}[theorem]{Definition}
\newtheorem{remark}[theorem]{Remark}
\newcommand{\E}{\mathbb E}
\newcommand{\Prob}{\mathbb P}
\newcommand{\Var}{\operatorname{Var}}
\newcommand{\sgn}{\operatorname{sgn}}
\newcommand{\atanh}{\operatorname{atanh}}
\newcommand{\He}{\operatorname{He}}
\newcommand{\dd}{\,\mathrm d}
\newcommand{\ii}{\mathrm i}
\newcommand{\norm}[1]{\left\lVert #1\right\rVert}
\newcommand{\abs}[1]{\left\lvert #1\right\rvert}
\newcommand{\crit}{\theta_{\mathrm c}}
\newcommand{\res}{\theta_{\mathrm r}}
\newcommand{\Disk}{\mathbb D}
\title{A sharp analytic renormalization threshold\texorpdfstring{\\}{ }for binary broadcasts on trees\texorpdfstring{\\[4pt]}{ }\large Critical-window residues, local asymptotics, and majority decoding}
\author{Research prepared for Vladimir Reshetnikov\\Developed with OpenAI ChatGPT}
\date{October 8, 2026}
\hypersetup{pdftitle={A sharp analytic renormalization threshold for binary broadcasts on trees},pdfauthor={Research prepared for Vladimir Reshetnikov; developed with OpenAI ChatGPT}}

\begin{document}
\maketitle
\begin{abstract}
Let $S_n$ be the sum of the spins at level $n$ of a rooted $b$-ary
binary broadcast, with the root fixed to $+1$ and edge correlation
$0<\theta\leq1$. Put $X_n=(b\theta)^{-n}S_n$. We prove that the
Gaussian-renormalized Laplace transforms
$\exp(-\Var(X_n)t^2/2)\E e^{tX_n}$ converge locally holomorphically
near zero if and only if $\theta>b^{-3/4}$. The limiting residue is
nonvanishing and depends analytically on $\theta$; convergence is
exponential, uniformly on compact parameter intervals above the threshold.
The converse follows from explicit third- and fourth-cumulant divergences,
including the boundary case. This fourth-cumulant threshold is distinct
from the classical reconstruction threshold $b^{-1/2}$.

For the finite-depth window
$\theta_n=b^{-1/2}\exp(\gamma/(2n))$, we identify a single analytic
residue controlling all bounded real $\gamma$, obtain expansions to
every order in $n^{-1}$, and prove a uniform lattice local limit expansion
to arbitrary order. A separate Fourier argument gives the full expansion
of majority-decoding success. On the binary tree at criticality its
$n^{-3/2}$ correction vanishes, and the next term is
$-5/(4\sqrt\pi\,n^{5/2})$. The proofs use exact broadcast recursions,
analytic contractions, and conditional independence; no factor-of-IID
claim is assumed. Exact rational computations verify the coefficient
identities and independent finite probability distributions. The classical
central limit theorem and earlier moment recursions are credited; the
precise novelty of the stronger theorem package remains subject to
specialist literature review.
\end{abstract}

\newpage
\tableofcontents
\newpage

\section{Scope, context, and the questions addressed}

Binary broadcasting separates two sources of randomness: a spin at the
root and independent errors along the edges. The exponentially many
leaves are strongly dependent because their ancestral paths overlap.
Nevertheless, the sum of the leaf spins admits an exact recursion.
This note asks how much information is retained by that recursion after
its growing Gaussian part has been removed.

There are three related questions. First, for which correlations does
Gaussian renormalization leave a genuine analytic limiting function?
Second, what happens when the correlation approaches the reconstruction
threshold at the same time as the tree grows? Third, can the analytic
information be converted into individual lattice probabilities and a
quantitative decoding calculation? The answers require control of the
whole transform, rather than only a fixed finite list of moments.

The principal result is a sharp threshold for locally holomorphic
convergence. Its value, $b^{-3/4}$, is forced by the fourth cumulant.
The familiar second-moment threshold is $b^{-1/2}$. Between these two
values the variance of the mean-normalized statistic diverges while its
Gaussian-renormalized transform converges analytically. At the second
threshold the variance grows linearly in the depth, and the residue gives
particularly explicit finite-size information.

\subsection{Relation to established work}

Limit theorems for multitype branching processes go back to Kesten and
Stigum~\cite{KestenStigum1966}. Moore and Snell~\cite{MooreSnell1979}
applied branching-process limit theorems to a phase transition for the
Ising model on a tree. The broadcast/reconstruction interpretation and
the role of the square of the edge correlation are developed in
Evans, Kenyon, Peres, and Schulman~\cite{EKPS2000}. General additive
functionals of bifurcating Markov chains also have subcritical, critical,
and supercritical fluctuation regimes; see Bitseki Penda and
Delmas~\cite{BitsekiDelmas2021}.

Exact recursions and moment asymptotics are themselves established tools.
M\'elin~\cite{Melin1996} derived recursions and the leading scaling of
all raw moments for the magnetization of the whole finite Cayley tree.
Our observable is the last level, and our main assertions concern a
locally holomorphic residue, its exact threshold of existence, a
finite-depth critical window, and uniform local expansions. We do not
claim to introduce the broadcast recursion or the critical central limit
theorem.

The general language of mod-Gaussian convergence was developed by
Jacod, Kowalski, and Nikeghbali~\cite{JKN2011}. M\'eliot and
Nikeghbali~\cite{MeliotNikeghbali2014} applied it to statistical-mechanical
models, including the one-dimensional Ising model. We use a locally
holomorphic version, stated explicitly below. Where the Gaussian
variance tends to a finite number, the same analytic renormalization
still makes sense, but it is not mod-Gaussian convergence with a
diverging variance parameter.

Gu, Roozbehani, and Polyanskiy~\cite{GuRoozbehaniPolyanskiy2020}
study full-boundary reconstruction near criticality. That question
involves all leaf spins, not only their sum, and a different order of
limits. Our majority calculation concerns a specified decoder. It is not
an assertion that majority is optimal among all boundary estimators.

\subsection{Repository connection and status of the contributions}

The motivating repository source is family 236 of the
\texttt{openai/math} collection: the manuscript on the sharp
factor-of-IID threshold for free Ising states on regular
trees~\cite{OpenAIFactor}. Its introduction explicitly separates
reconstruction, factor constructions, and finitary coding. We use its
topic as a research starting point. Its factor theorem, including its
critical case, is not an input to any proof in this note.

The repository was inspected at commit
\texttt{adc7f1241b42e322a6451854ab7e4b4c146bf78a}. Its README describes
a collection with results at different stages of verification. Accordingly,
all results used here are either proved below or attributed to established
primary literature for context only. The proposed contribution is the
combined analytic threshold, parameter-uniform residue, critical-window
expansion, and lattice/decoder consequences. The derivations are
self-contained; a claim of priority for every component would require a
more exhaustive specialist bibliography than the source audit accompanying
this research session.

\section{Model, normalizations, and main statements}

Fix an integer $b\geq2$. On the rooted $b$-ary tree, fix the root spin
to $+1$. For every edge $e$ choose independently
\[
 \Prob(\varepsilon_e=1)=a:=\frac{1+\theta}{2},\qquad
 \Prob(\varepsilon_e=-1)=1-a,\qquad 0<\theta\leq1.
\]
The spin at a vertex is the product of the edge signs along its path
from the root. Define
\begin{equation}\label{eq:normalizations}
 S_{n,\theta}=\sum_{|v|=n}\sigma_v,
 \qquad B=b\theta,\qquad \mu=\sqrt b\,\theta,
 \qquad X_{n,\theta}=B^{-n}S_{n,\theta},
 \qquad W_{n,\theta}=b^{-n/2}S_{n,\theta}=\mu^nX_{n,\theta}.
\end{equation}
The mean normalization is $X$, and the square-root-of-level-size
normalization is $W$. All probabilities and expectations below condition
on the root spin being $+1$, unless explicitly stated otherwise.

\begin{definition}
Write $V_{n,\theta}=\Var(X_{n,\theta})$. We say that Gaussian
renormalization converges locally holomorphically if, for some $\rho>0$,
\begin{equation}\label{eq:renormalized-transform}
 \Psi_{n,\theta}(t)
   :=\exp\!\left(-\frac{V_{n,\theta}t^2}{2}\right)
      \E e^{tX_{n,\theta}}
\end{equation}
converges uniformly on compact subsets of $\Disk_\rho=\{t:|t|<\rho\}$.
The limit has value $1$ at zero. In particular it is nonzero on a
possibly smaller disk.
\end{definition}

\begin{theorem}[Sharp analytic threshold]\label{thm:sharp}
For every $b\geq2$ and $0<\theta\leq1$, convergence in the preceding
definition holds if and only if
\begin{equation}\label{eq:threshold}
                         \theta>\res:=b^{-3/4}.
\end{equation}
Above this threshold there is a unique analytic germ $h_\theta$ with
\[
 h_\theta(0)=0,\qquad h_\theta'(0)=1,
 \qquad h_\theta''(0)=0
\]
satisfying
\begin{equation}\label{eq:fixed-h}
 h_\theta(t)=b\log\!\left(
        a e^{h_\theta(t/B)}+(1-a)e^{h_\theta(-t/B)}\right)
                -\frac{1-\theta^2}{2b\theta^2}t^2.
\end{equation}
The limit in \eqref{eq:renormalized-transform} is
$\Psi_\theta=e^{h_\theta}$.

For each compact real interval $I\subset(b^{-3/4},1]$, there are
$\rho,C>0$ and $q\in(0,1)$ such that analytic logarithms
$h_{n,\theta}=\log\Psi_{n,\theta}$ can be chosen on $\Disk_\rho$ and
\begin{equation}\label{eq:uniform-exponential}
 |h_{n,\theta}(t)-h_\theta(t)|\leq Cq^n|t|^3
 \qquad(\theta\in I,\ |t|\leq\rho/2).
\end{equation}
The germ $h_\theta(t)$ is jointly holomorphic in $t$ and in a complex
neighborhood of every real parameter in this interval.
\end{theorem}

Uniqueness here is uniqueness of the normalized analytic germ, not a
claim about arbitrary global solutions of the functional equation.
The endpoint $\theta=1$ gives $X_{n,1}=1$ and $h_1(t)=t$.

\begin{theorem}[Finite-depth critical window]\label{thm:window}
Let $\crit=b^{-1/2}$ and, for bounded real $\gamma$, put
\begin{equation}\label{eq:window}
 \theta_n=\crit e^{\gamma/(2n)},\qquad
 m_\gamma=e^{\gamma/2},\qquad W_n=W_{n,\theta_n}.
\end{equation}
For sufficiently large $n$, $0<\theta_n<1$. Then
$\E W_n=m_\gamma$ and
\begin{equation}\label{eq:window-variance}
 v_n:=\Var(W_n)
 =\left(1-\frac{e^{\gamma/n}}b\right)
          \frac{e^\gamma-1}{e^{\gamma/n}-1},
 \qquad v_n\big|_{\gamma=0}=\frac{b-1}{b}n.
\end{equation}
Uniformly for $\gamma$ in any fixed compact real interval, on a common
complex disk around zero,
\begin{equation}\label{eq:window-residue}
 e^{-v_nt^2/2}\E e^{tW_n}
      =\Psi_{\theta_n}(m_\gamma t)+O(q^n)
      \longrightarrow\Psi_{\crit}(m_\gamma t).
\end{equation}
For every $J\geq0$, the left side has an expansion through order
$n^{-J}$, with uniform remainder $O(n^{-J-1})$; its coefficients are
given explicitly by parameter derivatives of $\Psi_\theta$ in
\eqref{eq:window-taylor} below.
\end{theorem}

The remaining main assertions are the arbitrary-order uniform lattice
local expansion in Theorem~\ref{thm:local} and the majority expansion in
Theorem~\ref{thm:majority}. Their proofs include the high-frequency bounds
needed to pass from a small complex disk to lattice probabilities.

\section{Exact recursions and the parity decomposition}

\begin{lemma}[Distributional recursion and variance]\label{lem:recursion}
With independent copies $X_{n,\theta}^{(1)},\ldots,X_{n,\theta}^{(b)}$
and independent edge signs $\varepsilon_1,\ldots,\varepsilon_b$,
\begin{equation}\label{eq:distributional}
 X_{n+1,\theta}\ \overset{d}{=}
       \frac1B\sum_{j=1}^b\varepsilon_jX_{n,\theta}^{(j)}.
\end{equation}
Consequently
\begin{align}
 \E X_{n,\theta}&=1,\label{eq:mean}\\
 V_{n+1,\theta}&=\frac{V_{n,\theta}+1-\theta^2}{b\theta^2},\label{eq:variance-rec}\\
 V_{n,\theta}&=(1-\theta^2)\sum_{k=1}^n\mu^{-2k},\label{eq:variance-X}\\
 \Var(W_{n,\theta})&=(1-\theta^2)\sum_{j=0}^{n-1}\mu^{2j}.\label{eq:variance-W}
\end{align}
\end{lemma}

\begin{proof}
Conditional on the spins at the children of the root, the descendant
subtrees are independent. The spin sum in the $j$th subtree is its
root-fixed positive version multiplied by $\varepsilon_j$. Dividing by
$B^{n+1}$ proves \eqref{eq:distributional}. If a variable $X$ has mean
$1$ and variance $V$, then
$\Var(\varepsilon X)=V+1-\theta^2$. Independence gives
\eqref{eq:variance-rec}, and iteration gives the remaining formulas.
\end{proof}

For an analytic logarithm near zero, define
\[
 h_{n,\theta}(t)=\log\E e^{tX_{n,\theta}}-\frac12V_{n,\theta}t^2.
\]
The normalizations force $h_n=t+O(t^3)$. From the distributional
recursion and the variance identity we obtain the exact autonomous map
\begin{equation}\label{eq:autonomous}
 h_{n+1}(t)=b\log\!\left(ae^{h_n(t/B)}+(1-a)e^{h_n(-t/B)}\right)
                      -\frac{1-\theta^2}{2b\theta^2}t^2,
 \qquad h_0(t)=t.
\end{equation}
The growing variance has disappeared from this equation. The fact that
$t^2$ is even is essential: the Gaussian factor is identical in the two
terms of the mixture.

Write
\[
 Q_n(t)=\frac{h_n(t)-h_n(-t)}2,\qquad
 R_n(t)=\frac{h_n(t)+h_n(-t)}2.
\]
Thus $Q_n$ is odd with $Q_n=t+O(t^3)$, and $R_n$ is even with
$R_n=O(t^4)$. Put
\begin{equation}\label{eq:FG}
 F_\theta(u)=b\atanh(\theta\tanh u),\qquad
 G_\theta(u)=\frac b2\log\!\left(1+(1-\theta^2)\sinh^2u\right),
 \qquad d_\theta=\frac{1-\theta^2}{2b\theta^2}.
\end{equation}
All branches are the ones analytic at zero. Taking the odd and even
parts of \eqref{eq:autonomous} gives
\begin{align}
 Q_{n+1}(t)&=F_\theta\bigl(Q_n(t/B)\bigr),\label{eq:Qrec}\\
 R_{n+1}(t)&=bR_n(t/B)+G_\theta\bigl(Q_n(t/B)\bigr)-d_\theta t^2.
 \label{eq:Rrec}
\end{align}
To verify these formulas directly, factor out $e^{R_n}$ and use
\[
 ae^u+(1-a)e^{-u}=\cosh u+\theta\sinh u.
\]
The logarithmic ratio of its values at $u$ and $-u$ is
$2\atanh(\theta\tanh u)$, and their product is
$1+(1-\theta^2)\sinh^2u$.

The quadratic coefficient of $G_\theta(Q_n(t/B))$ is exactly
$d_\theta t^2$. The even forcing in \eqref{eq:Rrec} therefore starts
at degree four for every $n$. This cancellation produces the threshold
in Theorem~\ref{thm:sharp}.

\section{Analytic convergence above the sharp threshold}

We give the uniform analytic estimates in detail. They also justify all
later coefficient differentiations.

\subsection{A contraction for the odd part}

For $\rho>0$ and an analytic function vanishing to order at least $j$,
write
\[
 \norm{f}_{j,\rho}=\sup_{0<|t|\leq\rho}\frac{|f(t)|}{|t|^j}.
\]
We use the Banach space of analytic functions on $\Disk_\rho$ continuous
on its closure, with this weighted supremum norm. Odd and even
subspaces are closed. If necessary, one first works on a slightly
larger disk and subsequently decreases $\rho$.

\begin{lemma}[Uniform odd contraction]\label{lem:odd-contraction}
Let $I$ be a compact interval of real parameters with $b\theta>1$ on
$I$. There is a common disk on which the iterates in
\eqref{eq:Qrec}, starting at $Q_0=t$, converge exponentially to a
unique normalized odd analytic germ $Q_\theta$. The convergence and
the disk are uniform for $\theta\in I$, and $Q_\theta$ is jointly
holomorphic in the parameter and the argument.
\end{lemma}

\begin{proof}
Enlarge $I$ to a sufficiently small complex neighborhood $U$ so that
$|B|>1$ and
\[
 A_0:=\sup_{\theta\in U}|B|^{-2}<1.
\]
The functions $F_\theta(u)$ are jointly holomorphic for $\theta\in U$
and $u$ in a common neighborhood of zero. Uniformly there,
\[
 F_\theta(0)=0,\qquad F_\theta'(0)=B,
 \qquad F_\theta(u)=Bu+O(u^3).
\]
Choose a number $a_0$ with $A_0<a_0<1$. On odd functions $q=t+f$,
consider
\[
             (\mathcal A_\theta q)(t)=F_\theta(q(t/B)).
\]
For $\norm{f}_{3,\rho}\leq M$, the argument $q(t/B)$ lies in a
common small disk if $\rho$ is sufficiently small. The segment between
any two such arguments lies in the same disk. By the integral form of
the mean-value identity for holomorphic functions,
\begin{align*}
 \norm{\mathcal A_\theta q-\mathcal A_\theta\widetilde q}_{3,\rho}
 &\leq \left(\sup_{|u|\leq C\rho}|F_\theta'(u)|\right)
            |B|^{-3}\norm{q-\widetilde q}_{3,\rho}\\
 &\leq a_0\norm{q-\widetilde q}_{3,\rho}.
\end{align*}
The last inequality follows from $F_\theta'(u)=B+O(u^2)$ uniformly,
by decreasing $\rho$. Also
\[
       \norm{\mathcal A_\theta(t)-t}_{3,\rho}\leq C_0
\]
uniformly on $U$. Choose $M>C_0/(1-a_0)$ before decreasing $\rho$.
Then the closed ball $\norm{q-t}_{3,\rho}\leq M$ is invariant and
$\mathcal A_\theta$ is a contraction there. Its iterates converge with
an error bounded by $Ca_0^n$ in the weighted norm.

The iterates are jointly holomorphic on $U\times\Disk_\rho$ and
converge uniformly on smaller closed products. Their limit is jointly
holomorphic. To see uniqueness as a germ without a prescribed ball,
compare Taylor coefficients: the coefficient of $t^j$, for odd $j\geq3$,
has multiplier $B^{1-j}$, different from $1$, and its remaining forcing
depends only on lower coefficients. Every normalized analytic germ
therefore has the same Taylor series as the constructed one.
\end{proof}

\subsection{The even response and the fourth-degree multiplier}

Define the limiting forcing
\begin{equation}\label{eq:J}
 J_\theta(t)=G_\theta\bigl(Q_\theta(t/B)\bigr)-d_\theta t^2.
\end{equation}
It is even and vanishes to order four. Its finite-stage analogue is
$J_{n,\theta}$ with $Q_\theta$ replaced by $Q_n$.

\begin{lemma}[Even convergence]\label{lem:even-contraction}
Suppose, in addition, that $b/|B|^4<1$ uniformly on the parameter
neighborhood. Then
\begin{equation}\label{eq:Rseries}
 R_\theta(t)=\sum_{j=0}^{\infty}b^jJ_\theta(t/B^j)
\end{equation}
converges normally on a common disk, and $R_n\to R_\theta$
exponentially there, uniformly in the parameter. It is the unique even
analytic germ of order at least four satisfying
$R_\theta(t)=bR_\theta(t/B)+J_\theta(t)$.
\end{lemma}

\begin{proof}
On functions of order at least four the linear map
$\mathcal L_\theta R(t)=bR(t/B)$ has norm at most
\[
                     A_1:=\sup_{\theta\in U}\frac b{|B|^4}<1.
\]
This proves normal convergence of \eqref{eq:Rseries}. It also shows
uniqueness in the weighted space. Uniqueness as a germ follows again
from the multipliers $bB^{-j}\neq1$ for even $j\geq4$.

Because $G_\theta'(u)=O(u)$ near zero, the odd convergence implies
\[
 \norm{J_{n,\theta}-J_\theta}_{4,\rho}\leq C a_0^n.
\]
Indeed, the argument difference is $O(a_0^n t^3)$ and multiplication
by the derivative of $G$ supplies one more factor $t$.
Writing $e_n=\norm{R_n-R_\theta}_{4,\rho}$ gives
\[
 e_{n+1}\leq A_1e_n+Ca_0^n.
\]
For any $q$ with $\max(A_1,a_0)<q<1$, the finite geometric convolution
is bounded by $C'q^n$. This also absorbs a possible factor $n$ when the
two underlying rates coincide. Joint holomorphy follows from normal
convergence of the series and of the iterates.
\end{proof}

\begin{proof}[Proof of sufficiency in Theorem~\ref{thm:sharp}]
If $\theta>b^{-3/4}$, then $B>b^{1/4}>1$ and $b/B^4<1$.
The two lemmas apply on a complex neighborhood of every compact real
interval above the threshold. Put $h_\theta=Q_\theta+R_\theta$.
Equations \eqref{eq:Qrec}--\eqref{eq:Rrec} give
\eqref{eq:fixed-h}, and the estimates give
\eqref{eq:uniform-exponential}.

It remains to ensure that the logarithms used in the argument really
represent the probability transforms on the common disk. Start with
$h_0=t$. The contraction construction keeps every odd argument in a
disk where $\cosh u+\theta\sinh u$ has no zero; its logarithm is the
branch equal to zero at $u=0$. Thus the recursively constructed
$h_n=Q_n+R_n$ is analytic on that disk. Exponentiating
\eqref{eq:autonomous} and restoring the known Gaussian factors yields
the exact moment-generating recursion. By induction it is precisely
$\Psi_{n,\theta}$. No uniform zero-free property of the unnormalized
moment-generating functions has been assumed. It is a conclusion of the
construction on this disk.

Uniqueness of $Q$ and $R$ proves uniqueness of $h$. When $\theta=1$,
$F_1(u)=bu$ locally, $G_1=0$, and the same formulas give $h_1=t$.
\end{proof}

\section{Cumulant obstruction and sharpness}

For $j\geq3$, let $\kappa_{j,n}$ be the $j$th cumulant of
$X_{n,\theta}$. Removing its variance changes none of these cumulants.
Set
\begin{equation}\label{eq:abbrev-cumulants}
 z=\theta^2,\qquad c=1-z,\qquad
 \alpha=B^{-2}=\frac1{b^2z},\qquad
 \beta=\frac b{B^4}=\frac1{b^3z^2}.
\end{equation}

\begin{lemma}[Third and fourth cumulants]\label{lem:cumulants}
Starting with $\kappa_{3,0}=\kappa_{4,0}=0$, one has
\begin{align}
 \kappa_{3,n+1}&=\alpha(\kappa_{3,n}-2c),\label{eq:k3rec}\\
 \kappa_{4,n+1}&=\beta\bigl(\kappa_{4,n}
                     +4c\kappa_{3,n}-2c(1-3z)\bigr).
                     \label{eq:k4rec}
\end{align}
If $\alpha<1$, then
\begin{equation}\label{eq:k3exact}
 \kappa_{3,n}=C_3(1-\alpha^n),\qquad
 C_3=-\frac{2c}{b^2z-1}.
\end{equation}
\end{lemma}

\begin{proof}
For a variable $X$ of mean $1$, independent of a sign
$\varepsilon$ of mean $\theta$, expansion of the first four raw
moments gives
\begin{align*}
 \kappa_3(\varepsilon X)&=\theta\kappa_3(X)-2\theta(1-\theta^2),\\
 \kappa_4(\varepsilon X)&=\kappa_4(X)
           +4(1-\theta^2)\kappa_3(X)-2+8\theta^2-6\theta^4.
\end{align*}
All terms containing $\Var(X)$ cancel. Cumulants add for independent
summands and scale homogeneously under multiplication by $B^{-1}$.
Equation \eqref{eq:distributional} now gives the recurrences, and the
first one is an elementary geometric sum.
\end{proof}

\begin{proposition}[Failure at and below $b^{-3/4}$]\label{prop:failure}
If $0<\theta\leq b^{-3/4}$, then either
$\kappa_{3,n}\to-\infty$ or $\kappa_{4,n}\to-\infty$.
Consequently the normalized transforms cannot converge locally
holomorphically near zero.
\end{proposition}

\begin{proof}
If $B\leq1$, then $\alpha\geq1$ and $c>0$. The formula
\[
             \kappa_{3,n}=-2c\sum_{j=1}^n\alpha^j
\]
already proves divergence.

Suppose next that $B>1$. Set
\[
 C_0=-2c(1-3z),\qquad
 N=C_0+4cC_3,\qquad D=-4cC_3>0.
\]
Then \eqref{eq:k4rec} becomes
\begin{equation}\label{eq:k4-forced}
 \kappa_{4,n+1}=\beta\bigl(\kappa_{4,n}+N+D\alpha^n\bigr).
\end{equation}
At the boundary $\theta=b^{-3/4}$ one has $\beta=1$,
$\alpha=b^{-1/2}<1$, and
$\kappa_{4,n}=nN+D(1-\alpha^n)/(1-\alpha)$. Moreover
\begin{equation}\label{eq:Nsign}
 N=-\frac{2(1-z)}{b^2z-1}
        \bigl(3+(b^2-1)z-3b^2z^2\bigr)<0.
\end{equation}
Here $z=b^{-3/2}$, so the expression in parentheses is at least
$3-3/b>0$. This proves linear negative divergence at the boundary.

Below that boundary but above $1/b$, we have $\beta>1>\alpha$.
Solving \eqref{eq:k4-forced} gives
\begin{equation}\label{eq:k4-growth}
 \lim_{n\to\infty}\beta^{-n}\kappa_{4,n}
       =\beta\left(\frac{N}{\beta-1}
                          +\frac{D}{\beta-\alpha}\right).
\end{equation}
The sign must be checked: $C_0$ can be positive for a small part of
the binary-tree parameter range. A direct simplification gives
\begin{align}
 N(\beta-\alpha)+D(\beta-1)
 &=C_0(\beta-\alpha)+4cC_3(1-\alpha)\notag\\
 &=-\frac{2c\alpha}{bz}
           \bigl(1+3(b-1)z-bz^2\bigr)<0.\label{eq:negative-amplitude}
\end{align}
The quadratic in parentheses is concave on $[0,1]$ and positive at
both endpoints: its values are $1$ and $2b-2$. All denominators in
\eqref{eq:k4-growth} are positive, so its limit is strictly negative.
This proves exponential negative divergence.

Finally, locally uniform convergence of $\Psi_{n,\theta}$ to an
analytic function with value $1$ at zero would imply local uniform
convergence of their normalized logarithms on a smaller disk. Cauchy's
formula would then force convergence of every fixed derivative at zero.
The cumulant divergences contradict this. Together with the preceding
section, this completes Theorem~\ref{thm:sharp}.
\end{proof}

\begin{remark}[What the converse does not say]
The obstruction concerns a complex neighborhood and hence all local
derivatives. It does not exclude every conceivable pointwise limit on
the imaginary axis, nor a different scaling or a higher-order
renormalization. Those are different questions.
\end{remark}

\section{The critical residue and computable coefficients}

At $\theta=\crit=b^{-1/2}$, $B=\sqrt b$ and $X_n=W_n$.
Write $c_b=(b-1)/b$. Then
\begin{equation}\label{eq:critical-mgf}
 \E e^{tW_n}=
      \exp\!\left(\frac12c_bnt^2+h_{n,\crit}(t)\right),
 \qquad h_{n,\crit}\longrightarrow h_b:=h_{\crit}.
\end{equation}
The odd and even limits satisfy
\begin{align}
 Q_b(\sqrt b\,t)&=b\atanh\bigl(b^{-1/2}\tanh Q_b(t)\bigr),
     &Q_b'(0)&=1,\label{eq:critical-Q}\\
 R_b(t)&=bR_b(t/\sqrt b)
   +\frac b2\log\!\left(1+c_b\sinh^2 Q_b(t/\sqrt b)\right)
   -\frac{c_b}2t^2.
 \label{eq:critical-R}
\end{align}
These equations, or the convergent series \eqref{eq:Rseries}, specify
the residue completely without a closed elementary expression.

Let
\begin{equation}\label{eq:coeff-def}
 k_j(b)=h_b^{(j)}(0),\qquad
 A_j(b)=\left.\frac{\mathrm d^j}{\mathrm dt^j}e^{h_b(t)}\right|_{t=0}.
\end{equation}
The $k_j$ for $j\geq3$ are limiting cumulants. The $A_j$ are derivatives
of a residue, not necessarily moments of a probability measure. In fact
$h_b''(0)=0$ while $h_b'''(0)\neq0$, so $e^{h_b}$ cannot be a genuine
moment-generating function of a random variable with mean $1$ and
variance zero.

\begin{proposition}[First coefficients and finite-depth corrections]
\label{prop:critical-coefficients}
The first coefficients are
\begin{align}
 k_1&=1,&k_2&=0,&k_3&=-\frac2b,&k_4&=-\frac{2(b+1)}{b^2},\label{eq:k34}\\
 k_5&=\frac{16}{b^2},&
 k_6&=\frac{8(b-2)(2b-1)}{b^3(b+1)}.\label{eq:k56}
\end{align}
In particular,
\begin{equation}\label{eq:Avalues}
 A_0=1,\quad A_1=1,\quad A_2=1,\quad
 A_3=1-\frac2b,\quad
 A_5=1-\frac{30}b+\frac6{b^2}.
\end{equation}
At finite depth,
\begin{align}
 \kappa_3(W_n)&=-\frac2b(1-b^{-n}),\label{eq:critical-k3n}\\
 \kappa_4(W_n)&=-\frac{2(b+1)}{b^2}(1-b^{-n})
                   +\frac{8(b-1)}{b^2}n b^{-n}.\label{eq:critical-k4n}
\end{align}
\end{proposition}

\begin{proof}
Substituting $z=1/b$ in Lemma~\ref{lem:cumulants} gives
\eqref{eq:k34} and the finite formulas. The $n b^{-n}$ term is the
resonance between the third-cumulant convergence rate and the
fourth-degree multiplier.

For completeness, coefficient extraction from \eqref{eq:critical-Q}
first gives
\[
       Q_b(t)=t-\frac{t^3}{3b}+\frac{2t^5}{15b^2}+O(t^7).
\]
Writing $u=t/\sqrt b$, expanding
$\log(1+c_b\sinh^2 Q_b(u))$ through degree six in
\eqref{eq:critical-R}, and dividing the coefficient of $t^{2j}$
by $1-b^{1-j}$ gives
\[
 R_b(t)=-\frac{b+1}{12b^2}t^4
       +\frac{(b-2)(2b-1)}{90b^3(b+1)}t^6+O(t^8).
\]
Multiplication by the appropriate factorials proves
\eqref{eq:k56}. Finally $A_3=1+k_3$ and
$A_5=1+10k_3+5k_4+k_5$ follow by exponentiating $h_b$.
\end{proof}

\begin{table}[ht]
\centering
\begin{tabular}{@{}rrrrrr@{}}
\toprule
$b$&$k_3$&$k_4$&$k_5$&$k_6$&$A_5$\\
\midrule
2&$-1$&$-3/2$&$4$&$0$&$-25/2$\\
3&$-2/3$&$-8/9$&$16/9$&$10/27$&$-25/3$\\
4&$-1/2$&$-5/8$&$1$&$7/20$&$-49/8$\\
5&$-2/5$&$-12/25$&$16/25$&$36/125$&$-119/25$\\
\bottomrule
\end{tabular}
\caption{Exact coefficients of the critical residue. The supplied script
computes through degree ten and verifies the truncated fixed equation.}
\label{tab:coefficients}
\end{table}

\subsection{A coefficient algorithm with rational arithmetic}

At criticality every coefficient of $Q_b$ and $R_b$ is rational when
$b$ is an integer, despite the occurrence of $\sqrt b$ in the recursion.
For a formal odd series $Q$, form
\[
 \Phi(Q)=\sum_{j\geq0}\frac{b^{-j}}{2j+1}\tanh(Q)^{2j+1},
 \qquad L(Q)=\log\!\left(1+c_b\sinh^2Q\right).
\]
For odd $r$, the new coefficient of $Q$ is
$b^{(1-r)/2}[t^r]\Phi(Q)$. For even $r\geq4$, the new coefficient
of $R$ is
$b^{1-r/2}([t^r]R+[t^r]L(Q)/2)$. The exponents of $b$ are integers.

To compute the fixed series through any finite degree, first solve the
odd coefficients in increasing degree. At odd degree $r\geq3$, set
the unknown coefficient temporarily to zero, compute its forcing, and
divide by $1-b^{(1-r)/2}$. Next compute each even coefficient by
dividing its forcing by $1-b^{1-r/2}$. These denominators are nonzero.
Formal exponentiation then gives the $A_j$. This is a finite exact
algorithm; analytic convergence follows separately from the preceding
proof, not from the existence of the formal recursion alone.

\section{Uniform asymptotics in the critical window}

\begin{proof}[Proof of Theorem~\ref{thm:window}]
In the window \eqref{eq:window}, $\mu_n=e^{\gamma/(2n)}$, so
$\mu_n^n=m_\gamma$. The mean and variance follow from
Lemma~\ref{lem:recursion}. For every bound $|\gamma|\leq G$, the
parameters $\theta_n$ eventually lie in a common compact interval
$I\subset(b^{-3/4},1)$. The arguments $m_\gamma t$ lie in the common
analytic disk when $t$ lies in a sufficiently small disk depending on
$G$. Equation \eqref{eq:uniform-exponential} gives
\[
 \log\E e^{tW_n}-\frac12v_nt^2
       =h_{\theta_n}(m_\gamma t)+O(q^n|t|^3).
\]
Exponentiating proves \eqref{eq:window-residue}.

For the expansion, let $D_\theta=\theta\partial_\theta$. Joint
holomorphy and Taylor's theorem applied to
$s\mapsto\Psi_{\crit e^s}(m_\gamma t)$ give
\begin{equation}\label{eq:window-taylor}
 e^{-v_nt^2/2}\E e^{tW_n}
 =\sum_{j=0}^{J}\frac1{j!}
       \left(\frac{\gamma}{2n}\right)^j
       \left[D_\theta^j\Psi_\theta(m_\gamma t)\right]_{\theta=\crit}
       +O(n^{-J-1}).
\end{equation}
Cauchy bounds on a common parameter disk make the remainder uniform
for bounded $\gamma$ and for $t$ in a smaller closed disk. The
exponentially small finite-iteration error is absorbed into that
remainder. The same statement holds for the logarithm with $\Psi$
replaced by $h$.
\end{proof}

The leading variance is
\begin{equation}\label{eq:variance-leading}
 \frac{v_n}{n}\longrightarrow
        c_b\frac{e^\gamma-1}{\gamma},
 \qquad \frac{e^\gamma-1}{\gamma}\Big|_{\gamma=0}=1.
\end{equation}
This is positive and bounded away from zero on each compact real
$\gamma$-interval. The next terms, also uniform on such intervals, are
\begin{equation}\label{eq:variance-series}
 \begin{split}
 v_n={}&nc_b\frac{e^\gamma-1}{\gamma}
 -\frac{b+1}{2b}(e^\gamma-1)
 +\frac{c_b\gamma(e^\gamma-1)}{12n}\\
 &-\frac{c_b\gamma^3(e^\gamma-1)}{720n^3}+O(n^{-5}).
 \end{split}
\end{equation}
They follow from the Taylor series of $(e^s-1)^{-1}$ after its simple
pole has been separated. The expressions are interpreted continuously
at $\gamma=0$.

\begin{corollary}[Cumulants throughout the window]\label{cor:window-cumulants}
For every fixed $r\geq3$,
\[
 \kappa_r(W_n)=m_\gamma^r h_{\theta_n}^{(r)}(0)+O(q^n)
       \longrightarrow e^{r\gamma/2}k_r(b),
\]
uniformly for bounded real $\gamma$. Every fixed cumulant has an
expansion to all orders in $n^{-1}$. Moreover
\[
                   \frac{W_n-m_\gamma}{\sqrt{v_n}}
                       \ \Longrightarrow\ N(0,1).
\]
\end{corollary}

\begin{proof}
Cauchy's formula applied on a fixed smaller disk gives the cumulant
assertions. For the central limit theorem, substitute
$t=\ii u/\sqrt{v_n}$ in the logarithmic transform and subtract the
mean term. The quadratic term is $-u^2/2$ and the remaining terms tend
to zero for every fixed real $u$.
\end{proof}

The window theorem is stronger than this central limit corollary.
It retains the non-Gaussian residue before the final $\sqrt n$
normalization, including every higher cumulant and its parameter
correction. Its uniformity also permits subsequent local Fourier
analysis without choosing a different neighborhood for each $\gamma$.

\section{High frequencies and an all-orders lattice local limit}

A locally analytic transform expansion alone does not prove a lattice
local limit theorem. We now establish the required bounds across an
entire fundamental Fourier interval.

Fix $G<\infty$, assume $|\gamma|\leq G$, and use the window
\eqref{eq:window}. Write
\[
 M_{k;\theta}(t)=\E e^{tW_{k,\theta}},\qquad
 \Delta_n=2b^{-n/2},\qquad T_n=\frac{\pi}{\Delta_n}.
\]
For $n\geq1$, the lattice containing $W_n$ is
\begin{equation}\label{eq:lattice}
 \mathcal L_n=\Delta_n(\mathbb Z+\delta_b),\qquad
 \delta_b=\begin{cases}0,&b\text{ even},\\[1pt]1/2,&b\text{ odd}.
 \end{cases}
\end{equation}
Only finitely many of its points have positive probability. Probabilities
at the remaining points are defined to be zero.

\begin{lemma}[Uniform Fourier control]\label{lem:fourier}
There are constants $c,r>0$ and $n_0$, depending on $b,G$, such that,
for $n\geq n_0$,
\begin{align}
 |M_{n;\theta_n}(\ii t)|&\leq e^{-cn t^2}
       &&\left(|t|\leq r b^{\lceil n/2\rceil/2}\right),
       \label{eq:fourier-mid}\\
 |M_{n;\theta_n}(\ii t)|&\leq e^{-ct^2}
       &&\left(|t|\leq T_n\right).
       \label{eq:fourier-full}
\end{align}
In particular, the contribution from any fixed neighborhood of zero
removed from $[-T_n,T_n]$ is exponentially small in $n$.
\end{lemma}

\begin{proof}
Condition first on all spins at level $n-1$. The $b^n$ leaf spins are
conditionally independent, and each has mean $\theta_n$ or
$-\theta_n$. The modulus of a single conditional characteristic
factor, at argument $s$, is
\[
       |\cos s\pm\ii\theta_n\sin s|
           =\sqrt{1-(1-\theta_n^2)\sin^2s}.
\]
Since $1-\theta_n^2$ is uniformly positive for large $n$,
\[
 |M_{n;\theta_n}(\ii t)|
 \leq\exp\!\left[-c_0b^n\sin^2(t/b^{n/2})\right].
\]
On $|t|\leq T_n=\pi b^{n/2}/2$, the inequality
$|\sin s|\geq2|s|/\pi$ for $|s|\leq\pi/2$ proves
\eqref{eq:fourier-full}.

For the stronger bound, let $m=\lceil n/2\rceil$ and $k=n-m$.
Conditional on the spins at level $m$, the $b^m$ descendant blocks
are independent. Each block contributes $b^{-m/2}$ times a copy of
$W_{k,\theta_n}$ or its negative. The corresponding characteristic
functions are complex conjugates, hence have equal moduli. Therefore
\begin{equation}\label{eq:block-fourier}
 |M_{n;\theta_n}(\ii t)|
      \leq\left|M_{k;\theta_n}(\ii t/b^{m/2})\right|^{b^m}.
\end{equation}
For every $k\leq n$, the factor
$\mu_n^k=e^{\gamma k/(2n)}$ is bounded above and away from zero,
and \eqref{eq:variance-W} gives
$\Var(W_{k,\theta_n})\geq c_1k$.
The analytic construction above supplies uniform bounds on the even
parts $R_{k,\theta_n}=O(t^4)$ for all $k$. The odd parts are purely
imaginary on the imaginary axis because their Taylor coefficients are
real. Thus, for $|u|$ in a sufficiently small common interval,
\[
 \log|M_{k;\theta_n}(\ii u)|
    =-\frac12\Var(W_{k,\theta_n})u^2
        +\operatorname{Re}h_{k,\theta_n}(\ii\mu_n^ku)
    \leq-\frac{c_1}2ku^2+C u^4
    \leq-c_2ku^2.
\]
Choose the interval small enough for the final inequality for all
$k\geq1$. Substitution into \eqref{eq:block-fourier} proves
\eqref{eq:fourier-mid}, since $k\geq n/3$ for large $n$.

The interval from a fixed positive radius to $r b^{m/2}$ has an
$e^{-cn t^2}$ bound. Beyond $r b^{m/2}$, the $e^{-ct^2}$ bound has
a lower endpoint exponential in $n$, so its integral is
superexponentially small. This proves the last assertion as well.
\end{proof}

Define the parameter-dependent residue derivatives
\begin{equation}\label{eq:parameter-A}
 A_j(\theta)=\Psi_\theta^{(j)}(0).
\end{equation}
Then $A_j(\crit)=A_j(b)$ from \eqref{eq:coeff-def}. Let
$\varphi(y)=(2\pi)^{-1/2}e^{-y^2/2}$, and use the probabilists'
Hermite polynomials
\[
 \He_j(y)=(-1)^je^{y^2/2}\frac{\mathrm d^j}{\mathrm dy^j}e^{-y^2/2}.
\]

\begin{theorem}[Uniform local expansion to arbitrary order]\label{thm:local}
For every integer $J\geq0$, uniformly for $|\gamma|\leq G$ and
$x\in\mathcal L_n$, with $y=x/\sqrt{v_n}$,
\begin{equation}\label{eq:local}
 \frac{\Prob(W_n=x)}{\Delta_n}
   =\frac{\varphi(y)}{\sqrt{v_n}}
      \sum_{j=0}^{J}
       \frac{m_\gamma^jA_j(\theta_n)}{j!v_n^{j/2}}\He_j(y)
      +O\!\left(v_n^{-(J+2)/2}\right).
\end{equation}
The error is an absolute uniform error for the probability divided by
the lattice mesh. It is not a relative tail estimate.
\end{theorem}

\begin{proof}
The offset lattice inversion formula is
\begin{equation}\label{eq:inversion}
 \frac{\Prob(W_n=x)}{\Delta_n}
    =\frac1{2\pi}\int_{-T_n}^{T_n}
          M_{n;\theta_n}(\ii t)e^{-\ii tx}\dd t.
\end{equation}
It holds also for the half-integer offset: after multiplication by
$e^{-\ii tx}$, all differences of support points are integer multiples
of $\Delta_n$.

Fix a small radius $r_1$ in the common analytic disk. Lemma~\ref{lem:fourier}
makes the integral over $r_1\leq|t|\leq T_n$ exponentially small.
On $|t|\leq r_1$, the uniform analytic expansion gives
\[
 M_{n;\theta_n}(\ii t)
   =e^{-v_nt^2/2}
       \left[\Psi_{\theta_n}(\ii m_\gamma t)+O(q^n|t|^3)\right].
\]
The residue has, uniformly in the parameters, the Taylor expansion
\[
 \Psi_{\theta_n}(\ii m_\gamma t)
   =\sum_{j=0}^{J}\frac{A_j(\theta_n)}{j!}(\ii m_\gamma t)^j
                       +O(|t|^{J+1}).
\]
The integral of the Taylor remainder is bounded by
$C\int_{\mathbb R}e^{-v_nt^2/2}|t|^{J+1}\dd t
=O(v_n^{-(J+2)/2})$. The finite-iteration error is exponential.
Extending the Gaussian-polynomial integrals to the real line changes
them by an exponentially small amount. Finally,
\[
 \frac1{2\pi}\int_{\mathbb R}e^{-vt^2/2}(\ii t)^j e^{-\ii tx}\dd t
       =v^{-(j+1)/2}\He_j(x/\sqrt v)\varphi(x/\sqrt v).
\]
All estimates used $|e^{-\ii tx}|=1$, so they are uniform in $x$.
\end{proof}

The coefficients in \eqref{eq:local} can themselves be expanded in
$n^{-1}$ by \eqref{eq:window-taylor}, and $v_n$ by
\eqref{eq:variance-series}. Keeping the exact variance and the exact
analytic parameter $\theta_n$ is usually a cleaner numerical formula.
The first two nonconstant terms are
\[
 \frac{m_\gamma}{\sqrt{v_n}}\He_1(y)
 +\frac{m_\gamma^2}{2v_n}\He_2(y),
\]
because $A_1(\theta)=A_2(\theta)=1$ throughout the parameter range.

\section{Majority decoding and a binary-tree cancellation}

Give the root a fair random sign for the decoding experiment, and
predict it by the sign of the sum of the leaves, breaking ties by an
independent fair coin. Spin-flip symmetry shows that its success
probability is
\begin{equation}\label{eq:majority-def}
 p_n=\Prob(W_n>0\mid\sigma_{\varnothing}=+1)
          +\frac12\Prob(W_n=0\mid\sigma_{\varnothing}=+1)
      =\frac12+\frac12\E[\sgn W_n\mid\sigma_{\varnothing}=+1].
\end{equation}
The conditioning agrees with all preceding definitions.

Integrating the absolute error in the local expansion over an unbounded
half-line would not be sufficient. We instead use an exact lattice
Fourier formula for the sign function.

\begin{lemma}[Lattice sign formula]\label{lem:sign}
Let $Y$ be finitely supported on $\Delta(\mathbb Z+\delta)$, with
$\delta=0$ or $1/2$, and let $M(t)=\E e^{tY}$. Then
\begin{equation}\label{eq:sign-integral}
 \E\sgn Y=\frac{\Delta}{\pi}
       \int_0^{\pi/\Delta}\operatorname{Im}M(\ii t)
           K_\delta(\Delta t/2)\dd t,
 \qquad K_0(u)=\cot u,\quad K_{1/2}(u)=\csc u.
\end{equation}
The apparent singularity at zero is removable after multiplication by
$\operatorname{Im}M(\ii t)$.
\end{lemma}

\begin{proof}
For an integer $k$,
\[
 \frac1\pi\int_0^\pi\sin(ks)\cot(s/2)\dd s=\sgn k.
\]
For a half-integer $k+1/2$,
\[
 \frac1\pi\int_0^\pi\sin((k+1/2)s)\csc(s/2)\dd s
                         =\sgn(k+1/2).
\]
For positive arguments these identities follow by expanding the sine
ratios as finite cosine sums; the constant term is $1$, and all
nonconstant integer cosine terms integrate to zero. Negative arguments
follow by oddness. Multiplying by the finitely many probabilities,
summing, and changing variables $s=\Delta t$ gives the formula.
\end{proof}

\begin{theorem}[All-orders majority expansion]\label{thm:majority}
For every integer $K\geq0$, uniformly for bounded real $\gamma$,
\begin{equation}\label{eq:majority-series}
 p_n=\frac12+\frac1{\sqrt{2\pi}}
      \sum_{j=0}^{K}
       \frac{(-1)^j m_\gamma^{2j+1}A_{2j+1}(\theta_n)}
            {(2j+1)2^j j!\,v_n^{j+1/2}}
       +O(v_n^{-K-3/2}).
\end{equation}
At exact criticality, $v_n=c_bn$ and
\begin{equation}\label{eq:majority-critical}
 p_n=\frac12+\frac1{\sqrt{2\pi}}
 \left\{v_n^{-1/2}
   -\frac{1-2/b}{6v_n^{3/2}}
   +\frac{1-30/b+6/b^2}{40v_n^{5/2}}
   +O(v_n^{-7/2})\right\}.
\end{equation}
In particular, for $b=2$,
\begin{equation}\label{eq:binary-cancellation}
 p_n=\frac12+\frac1{\sqrt{\pi n}}
                   -\frac5{4\sqrt\pi\,n^{5/2}}+O(n^{-7/2}).
\end{equation}
\end{theorem}

\begin{proof}
Apply Lemma~\ref{lem:sign} with $\Delta=\Delta_n$ and the correct
offset from \eqref{eq:lattice}. On $0<t\leq T_n$,
\[
 |\Delta_n K_{\delta_b}(\Delta_nt/2)|\leq C/t.
\]
Thus Lemma~\ref{lem:fourier} makes the contribution from $t\geq r_1$
exponentially small, for a fixed small $r_1>0$. On the remaining
interval,
\[
 \Delta_n K_{\delta_b}(\Delta_nt/2)=\frac2t+O(\Delta_n^2t),
\]
uniformly for both parity cases. The kernel replacement error is
exponentially small because $\Delta_n^2=4b^{-n}$ and the transform
has Gaussian decay there. Consequently
\[
 p_n-\frac12=\frac1\pi\int_0^{r_1}
       e^{-v_nt^2/2}\frac{\operatorname{Im}
                 \Psi_{\theta_n}(\ii m_\gamma t)}t\dd t
                      +O(q_1^n)
\]
for some $q_1<1$. Near zero the finite-iteration error is
$O(q^n t^3)$, so division by $t$ causes no problem.

All residue Taylor coefficients are real. Its imaginary part is
therefore an odd power series. Expanding it through degree $2K+1$
leaves $O(t^{2K+3})$. After division by $t$ and Gaussian integration,
the remainder is $O(v_n^{-K-3/2})$. For $j\geq0$,
\[
 \int_0^\infty t^{2j}e^{-v_nt^2/2}\dd t
       =\sqrt{\frac\pi2}(2j-1)!!\,v_n^{-j-1/2},
 \qquad (-1)!!=1.
\]
Using
$(2j-1)!!/(2j+1)!=1/((2j+1)2^j j!)$
gives \eqref{eq:majority-series}. The critical coefficients in
\eqref{eq:Avalues} give \eqref{eq:majority-critical}. For $b=2$,
$A_3=0$, $A_5=-25/2$, and $v_n=n/2$, which yields
\eqref{eq:binary-cancellation}.
\end{proof}

The vanishing term for the binary tree is a cancellation between the
mean shift and the limiting third cumulant. Keeping only a shifted
Gaussian would predict a nonzero correction of order $n^{-3/2}$.
The all-orders calculation detects why that prediction fails.

The leading window advantage is
\[
 p_n-\frac12\sim
 \frac{e^{\gamma/2}}{\sqrt{2\pi n c_b}}
          \sqrt{\frac{\gamma}{e^\gamma-1}},
\]
with continuous interpretation at $\gamma=0$. This concerns the
specified majority statistic and finite-depth window only.

\section{Finite even-cumulant renormalization beyond the threshold}

The parity decomposition gives a further consequence. A failed
Gaussian renormalization need not preclude renormalization by a finite
number of even cumulants.

\begin{theorem}[Higher even renormalization]\label{thm:higher}
Let $\theta>1/b$ and choose an integer $K\geq1$ such that
\begin{equation}\label{eq:higher-condition}
                         (b\theta)^{2K+2}>b.
\end{equation}
If $\kappa_{2j,n}$ denotes the actual even cumulant of
$X_{n,\theta}$, then
\begin{equation}\label{eq:higher-renorm}
 \exp\!\left(-\sum_{j=1}^{K}
             \frac{\kappa_{2j,n}}{(2j)!}t^{2j}\right)
                  \E e^{tX_{n,\theta}}
\end{equation}
converges locally holomorphically to a nonzero analytic function.
Convergence is exponential and locally uniform in $\theta$ whenever
\eqref{eq:higher-condition} holds uniformly. Thus every $\theta>1/b$
admits a finite even-cumulant renormalization of this form.
\end{theorem}

\begin{proof}
The odd logarithm $Q_n$ is unaffected by subtraction of any even
polynomial and converges by Lemma~\ref{lem:odd-contraction}, which
requires only $b\theta>1$. Let $H_n$ be the even part of the full
logarithm $\log\E e^{tX_n}$, and let $\Pi_{\leq2K}$ extract its
Taylor polynomial of degree at most $2K$. Set
$E_n=H_n-\Pi_{\leq2K}H_n$. Taking even parts before renormalization
gives
\[
 E_{n+1}(t)=bE_n(t/B)
    +(1-\Pi_{\leq2K})G_\theta(Q_n(t/B)).
\]
The linear map has norm $b/|B|^{2K+2}<1$ on even functions of order
at least $2K+2$. The forcing converges exponentially in that weighted
norm on a smaller disk. To justify the last statement, use the
uniform odd convergence on a larger disk and Cauchy's coefficient
bounds: after removing finitely many Taylor terms, the remaining tail
is bounded by $Cq^n|t|^{2K+2}$ on the smaller disk. The same geometric
convolution argument as in Lemma~\ref{lem:even-contraction} proves
convergence of $E_n$. The logarithm of \eqref{eq:higher-renorm} is
$Q_n+E_n$, proving the theorem.
\end{proof}

This theorem is a sufficient criterion for a specified finite
renormalization. It does not assert that every higher-order boundary
is sharp or that all possible cancellations have been classified.
For $K=1$, Theorem~\ref{thm:sharp} supplies the matching converse.
For $\theta\leq1/b$, the divergent third cumulant prevents convergence
after subtracting only even cumulants.

\section{Exact verification, reproducibility, and practical use}

The analytic arguments are mathematical proofs. The supplementary
computations serve two more limited purposes: they verify the finite
coefficient algebra and compare the transform recursion with independently
constructed finite distributions.

\subsection{What the supplied script checks}

The file \texttt{verify.py} uses only Python's standard library, with
\texttt{fractions.Fraction} for every assertion. It computes truncated
power series through degree ten and checks the following.
\begin{enumerate}[leftmargin=*,itemsep=3pt]
\item For $b=2,3,4,5,9$, the rational fixed series satisfies the
      truncated odd and even fixed equations exactly.
\item The limiting third and fourth cumulants agree with
      \eqref{eq:k34}, and $A_3=1-2/b$.
\item For these five branching values and every $0\leq n\leq20$,
      the finite third and fourth cumulants agree with
      \eqref{eq:critical-k3n}--\eqref{eq:critical-k4n}.
\item For $b=4$ through depth three and $b=9$ through depth two,
      an independent probability-generating recursion constructs the
      full law of the number of positive leaves. All raw moments
      through degree ten, after logarithmic conversion, agree with
      the parity recursion.
\end{enumerate}
The last checks use rational critical correlations $1/2$ and $1/3$,
respectively. They compare distributions with up to $81$ leaves, rather
than merely comparing two copies of the same cumulant formula.
Additional exact checks cover rational correlations above and below
$b\theta=1$, and verify the negative fourth-cumulant growth amplitude,
including the binary-tree case with positive initial fourth forcing.
The recorded execution reports \textbf{392 exact scalar checks} and
passes all fixed-series assertions. The JSON output retains exact
fractions, finite cumulants, and finite majority probabilities.

To reproduce the algebraic checks and the PDF, run from this directory:
\begin{verbatim}
python verify.py --output verification.json
latexmk -pdf -interaction=nonstopmode -halt-on-error article.tex
\end{verbatim}
The PDF requires standard LaTeX packages; the algebraic verification
requires no third-party Python packages. These calculations do not
verify the analytic contraction theorem by numerical sampling, and no
Lean formalization is included.

\subsection{How to use the expansions responsibly}

For a fixed desired expansion order, compute the critical residue
coefficients once by the triangular rational algorithm. At exact
criticality, the variance is known exactly. In the window, keeping
\eqref{eq:window-variance} unexpanded avoids unnecessary cancellation,
and the parameter-dependent residue can be iterated to exponential
accuracy inside its analytic disk. Parameter derivatives in
\eqref{eq:window-taylor} can also be computed from the differentiated
fixed equations.

The local theorem controls absolute lattice probabilities, uniformly
over the lattice. It does not yield a relative approximation deep in
the tails, where the Gaussian main term can be much smaller than the
absolute error. The majority theorem avoids accumulating this error
by using a separate exact Fourier identity. None of these formulas
supplies a computationally efficient sampler for the entire boundary
spin configuration, or a factor-of-IID coding radius.

\section{Further research questions}

The following are proposed research directions, not claims that each
formulation has been established as an open problem in the literature.

\begin{enumerate}[leftmargin=*,itemsep=6pt]
\item \textbf{Global domain of the residue.}
Determine the maximal complex continuation domain of $h_\theta$ and
the location and type of its first singularities. The local contraction
does not describe the zeros of finite-depth transforms far from zero.
The one-dimensional iteration in \eqref{eq:critical-Q} provides a
concrete entry point through analytic iteration theory.

\item \textbf{Sharp exponential convergence rates.}
Classify resonances among the odd multipliers $B^{1-j}$ and the even
multipliers $bB^{-j}$. The explicit $n b^{-n}$ correction in the
fourth cumulant suggests a full expansion in exponential scales times
polynomials in $n$. Such a result could strengthen the deliberately
coarse uniform $O(q^n)$ bound.

\item \textbf{Higher renormalization boundaries.}
For fixed $K$, decide whether \eqref{eq:higher-condition} is generically
or universally sharp after removal of cumulants through degree $2K$.
Identify exceptional parameters where the first obstructing forcing
vanishes. The fourth-cumulant sign calculation provides a model for
the required noncancellation arguments.

\item \textbf{Other tree observables.}
Develop the corresponding residue theory for the sum over all levels,
for a fixed-width band of levels, and for weighted level sums. This
would permit direct comparison with the total-magnetization
asymptotics of M\'elin. The added root contribution changes the
renormalization equation and must be retained explicitly.

\item \textbf{Irregular and multitype trees.}
Extend the analytic contraction to periodic branching, multitype
deterministic trees, and suitable random trees. In these settings the
scalar fourth-degree multiplier should be replaced by an operator,
and the correct normalization may be vector-valued.

\item \textbf{General channels.}
For finite-state broadcasts, identify when the leading Gaussian
covariance is invariant under the channel symmetries. The exact
parity cancellation used here need not survive. A useful first target
is a finite group channel with a real one-dimensional character.

\item \textbf{Information retained by the sum.}
Compute asymptotic mutual information between the root and the level
sum, with error estimates compatible with small probabilities. Compare
the optimal estimator based only on the sum with majority. The
present theorem evaluates majority and does not prove their equality.

\item \textbf{Full-boundary reconstruction in the depth window.}
Compare the residue-based majority constants with the critical-window
behavior of the complete posterior root magnetization. This requires
the nonlinear posterior recursion and cannot be inferred from a local
limit theorem for the sum.

\item \textbf{Certified computation and formalization.}
Formalize the finite generating-function and cumulant recurrences,
then isolate reusable analytic lemmas for parameter-uniform
contractions and lattice Fourier inversion. Interval bounds for the
contraction disk would turn the asymptotic residue into a certified
numerical object with explicit truncation error.
\end{enumerate}

\section{Conclusion and precise contribution boundary}

The exact broadcast recursion admits a useful separation: its odd
logarithm evolves autonomously, while its even logarithm is driven
linearly by the odd part. After the variance is removed, the first
remaining even degree is four. This gives the analytic threshold
$b^{-3/4}$, and an explicit cumulant calculation proves the threshold
is sharp, including the resonant endpoint.

Parameter-uniform analyticity then organizes the entire finite-depth
window around $b^{-1/2}$. Conditional independence at two different
cuts of the tree supplies the missing Fourier control, so the
analytic residue produces uniform lattice local expansions and an
independent all-orders majority calculation. These deductions require
neither a factor construction nor a newly assumed reconstruction
theorem. They provide a complete derivation of the stated results and
a reproducible foundation for further work, while leaving historical
priority and broader extensions to additional specialist review.

\appendix
\section{Source and claim audit}

The public repository source used to select the problem is
\begin{quote}\small
\path{preprints/The-sharp-factor-of-IID-threshold-for-the-free-Ising-model-on-regular-trees-September-26-2026/build/sections/01-introduction.tex}
\end{quote}
at the commit recorded in Section 1. Its role is motivational. The
current paper's theorem hypotheses are stated directly in terms of a
finite independent-edge broadcast and are checked within the paper.

\begin{longtable}{@{}>{\raggedright\arraybackslash}p{0.29\textwidth}>{\raggedright\arraybackslash}p{0.64\textwidth}@{}}
\toprule
Claim or method&Status in this note\\
\midrule
\endhead
Broadcast/Ising correspondence and classical fluctuation regimes&
Established background; credited to the primary literature.\\[4pt]
Exact broadcast recursion&
Elementary and established; derived explicitly for the chosen level-sum
normalization.\\[4pt]
Moment recursions and leading moment scaling&
Earlier work exists, including total-tree results; not claimed as a new
general method.\\[4pt]
Sharp local holomorphic residue threshold&
Proved in Theorem~\ref{thm:sharp}, with a sign-controlled converse.\\[4pt]
Parameter-uniform critical window&
Proved from a jointly analytic residue; exact variance and all orders
are supplied.\\[4pt]
Uniform lattice local expansion&
Proved with explicit conditioning bounds over the full Fourier interval.\\[4pt]
Majority-decoding expansion&
Proved by a separate lattice sign identity, including even/odd lattice
offsets and fair tie handling.\\[4pt]
Factor-of-IID theorem or finitary coding&
Neither proved nor assumed.\\[4pt]
Optimal full-boundary reconstruction&
Not claimed; the statistic and decoder are specified.\\[4pt]
Priority of every new formulation&
Not certified by the limited literature audit; independent derivation
and mathematical proof are the claims made here.\\
\bottomrule
\end{longtable}

\begin{thebibliography}{99}

\bibitem{KestenStigum1966}
H. Kesten and B. P. Stigum,
Additional limit theorems for indecomposable multidimensional
Galton--Watson processes,
\emph{Annals of Mathematical Statistics} \textbf{37} (1966), 1463--1481.
\url{https://doi.org/10.1214/aoms/1177699139}.

\bibitem{MooreSnell1979}
T. Moore and J. L. Snell,
A branching process showing a phase transition,
\emph{Journal of Applied Probability} \textbf{16} (1979), 252--260.
\url{https://doi.org/10.2307/3212894}.

\bibitem{EKPS2000}
W. Evans, C. Kenyon, Y. Peres, and L. J. Schulman,
Broadcasting on trees and the Ising model,
\emph{Annals of Applied Probability} \textbf{10} (2000), 410--433.
\url{https://doi.org/10.1214/aoap/1019487349}.

\bibitem{Melin1996}
R. M\'elin,
Magnetization and overlap distributions of the ferromagnetic Ising model
on the Cayley tree,
\emph{Journal de Physique I} \textbf{6} (1996), 1435--1450.
\url{https://doi.org/10.1051/jp1:1996156};
\url{https://arxiv.org/abs/cond-mat/9607215}.

\bibitem{BitsekiDelmas2021}
S. Val\`ere Bitseki Penda and J.-F. Delmas,
Central limit theorem for bifurcating Markov chains under point-wise
ergodic conditions, preprint, version 2 (2021).
\url{https://arxiv.org/abs/2012.04741v2}.

\bibitem{JKN2011}
J. Jacod, E. Kowalski, and A. Nikeghbali,
Mod-Gaussian convergence: new limit theorems in probability and number
theory, \emph{Forum Mathematicum} \textbf{23} (2011), 835--873.
\url{https://doi.org/10.1515/form.2011.030}.

\bibitem{MeliotNikeghbali2014}
P.-L. M\'eliot and A. Nikeghbali,
Mod-Gaussian convergence and its applications for models of statistical
mechanics, preprint (2014).
\url{https://arxiv.org/abs/1409.2849}.

\bibitem{GuRoozbehaniPolyanskiy2020}
Y. Gu, H. Roozbehani, and Y. Polyanskiy,
Broadcasting on trees near criticality,
\emph{2020 IEEE International Symposium on Information Theory}, 2020;
preprint \url{https://arxiv.org/abs/2005.07801}.

\bibitem{OpenAIFactor}
OpenAI,
The sharp factor-of-IID threshold for the free Ising model on regular
trees, repository manuscript, September 26, 2026.
\url{https://github.com/openai/math/tree/adc7f1241b42e322a6451854ab7e4b4c146bf78a/preprints/The-sharp-factor-of-IID-threshold-for-the-free-Ising-model-on-regular-trees-September-26-2026}.
Consulted as motivation; its main theorem is not assumed here.

\end{thebibliography}
\end{document}
